\documentclass[12pt]{article}
\usepackage[T1]{fontenc}
\usepackage[utf8]{inputenc}
\usepackage{lmodern}
\usepackage{textcomp}
\usepackage{enumitem}
\usepackage{geometry}
\usepackage{graphicx}
\usepackage{amsmath, amssymb}
\geometry{margin=1in}
\begin{document}
\begin{center}
Sociology - Undergraduate Level Assessment 5 \
Total Marks: 40 \
Answer all questions.
\end{center}

\section*{Multiple Choice (10 marks)}
\begin{enumerate}[label=\arabic*.]
\item In Durkheim's work, the term 'collective representations' refers to:
\begin{enumerate}[label=(\alph*)]
    \item effervescent ceremonies that create a feeling of belonging
    \item images of gods or totems that are widely recognized
    \item shared ideas and moral values, often symbolized by an object or figurehead
    \item ideological tools used to obscure class divisions
\end{enumerate}
\item Scientific management involved:
\begin{enumerate}[label=(\alph*)]
    \item the subdivision of labour into small tasks
    \item the measurement and specification of work tasks
    \item motivation and rewards for productivity
    \item all of the above
\end{enumerate}
\item There was a growth in income inequality in the 1980 s because:
\begin{enumerate}[label=(\alph*)]
    \item rates of income tax increased equally for all occupational groups
    \item there were more professional jobs available but not enough people to fill them
    \item the price of consumer goods rose at a higher rate than earnings
    \item growing unemployment made more people dependent on welfare benefits
\end{enumerate}
\item Post-Fordist forms of media production and consumption involve:
\begin{enumerate}[label=(\alph*)]
    \item the mass production of standardized products for passive audiences
    \item television based on producer-broadcaster rather than publisher-broadcaster models
    \item a diverse range of products aimed at niche markets
    \item increasing numbers of advertisements for motoring and car-related products
\end{enumerate}
\item Theories of racialized discourse suggest that:
\begin{enumerate}[label=(\alph*)]
    \item race is an objective way of categorizing people on biological grounds
    \item the idea of race is socially constructed through powerful ideologies
    \item race relations in Britain and America can be traced back to colonial times
    \item people choose their racial identity and this becomes fixed
\end{enumerate}
\end{enumerate}

\section*{True / False (10 marks)}
\textit{Answer True or False}
\begin{enumerate}[label=\arabic*.]
\setcounter{enumi}{5}
\item True or False: The social democratic model is characterized by loyalty to the state and adherence to traditional values.
\item True or False: Differential association theory posits that individuals are inherently predisposed to criminal behavior regardless of their social environment.
\item True or False: Definitions and indicators of social class can vary, making valid comparisons across studies problematic.
\item True or False: Goldthorpe identified the 'service class' as those in non-manual occupations, exercising authority on behalf of the state.
\item True or False: The concept of 'culture industry' emphasizes the emergence of subcultures and counter-cultures within society.
\end{enumerate}

\section*{Long Form Response (20 marks)}
\textit{Answer each question with a clear and complete explanation.}
\begin{enumerate}[label=\arabic*.]
\setcounter{enumi}{10}
\item Discuss the potential social implications of a new highway extending from a major city through farmland, particularly focusing on the phenomenon of suburbanization and its effects on local communities.
\item Discuss the societal changes that characterized the 'sexual revolution' of the 1960 s, and analyze the factors that contributed to the emerging fear of HIV and AIDS in later years.
\end{enumerate}

\vfill
\noindent\textit{End of Paper}
\end{document}